\section*{अध्याय \ref{ch:b1:faults-folds-fabrics} --- भ्रंश, वलन और विन्यास}

\begin{solution}{exo:b1:faults-folds-fabrics:1}
(a) एक \omterm{def:b1:faults-folds-fabrics:fault-types}{सामान्य भ्रंश}। (b) एक \omterm{def:b1:faults-folds-fabrics:fault-types}{उत्क्रम भ्रंश}, लगभग शुद्ध नति-सरकन। (c) एक \omterm{def:b1:faults-folds-fabrics:fault-types}{बाएँ-पार्श्वीय} \omterm{def:b1:faults-folds-fabrics:fault-types}{स्ट्राइक-स्लिप भ्रंश}। (d) एक \omterm{def:b1:faults-folds-fabrics:fault-types}{दाएँ-पार्श्वीय} \omterm{def:b1:faults-folds-fabrics:fault-types}{स्ट्राइक-स्लिप भ्रंश}। (e) एक \omterm{def:b1:faults-folds-fabrics:fault-types}{तिर्यक-सरकन भ्रंश}: \omterm{def:b1:faults-folds-fabrics:fault-elements}{लटकती दीवार} नीचे और \omterm{def:b1:geological-maps:attitude}{स्ट्राइक} के अनुदिश बाईं ओर खिसकी है; उसका स्ट्राइक-स्लिप घटक, \omterm{def:b1:faults-folds-fabrics:fault-elements}{निवल सरकन} का $\cos\ang{40} = 0.77$, उसके नति-सरकन घटक $0.64$ से बड़ा है, इसलिए उसे सामान्य घटक वाला एक \omterm{def:b1:faults-folds-fabrics:fault-types}{बाएँ-पार्श्वीय} भ्रंश कहा जाता है। (f) \ang{25} की \omterm{def:b1:geological-maps:attitude}{नति} पर, जो \ang{45} से कम है, \omterm{def:b1:faults-folds-fabrics:fault-types}{उत्क्रम भ्रंश} एक \omterm{def:b1:faults-folds-fabrics:fault-types}{अधिक्षेप भ्रंश} है।
\end{solution}

\begin{solution}{exo:b1:faults-folds-fabrics:2}
$|\sin r| = 1$ के साथ भ्रंश का \omterm{def:b1:cross-sections:throw}{विस्थापन} $300\sin\ang{55} = \qty{245.7}{m}$ है, और
\omterm{def:b1:cross-sections:throw}{क्षैतिज विस्थापन} $300\cos\ang{55} = \qty{172.1}{m}$।
\end{solution}

\begin{solution}{exo:b1:faults-folds-fabrics:3}
विपरीत दिशाओं में नत पार्श्वों वाले सीधे \omterm{def:b1:geological-maps:fold}{वलन} के लिए, जिसके \omterm{def:b1:faults-folds-fabrics:fold-elements}{पार्श्व} $\delta_1$ और $\delta_2$ नत हैं, $\iota = \ang{180} - \delta_1 - \delta_2$। (a) \ang{110}: खुला। (b) \ang{35}: सघन।
(c) \ang{7}: समनत।
\end{solution}

\begin{solution}{exo:b1:faults-folds-fabrics:4}
\omterm{def:b1:faults-folds-fabrics:fault-types}{सामान्य भ्रंश} $\ang{45} + \ang{35}/2 = \ang{62.5}$ पर नत हैं, अधिक्षेप $\ang{45} - \ang{17.5} = \ang{27.5}$ पर। संयुग्मी \omterm{def:b1:faults-folds-fabrics:fault-types}{स्ट्राइक-स्लिप भ्रंश} प्रत्येक $\sigma_1$ के साथ \ang{27.5} का कोण बनाते हैं, उसके दोनों ओर, इसलिए वे \ang{55} पर मिलते हैं।
\end{solution}

\begin{solution}{exo:b1:faults-folds-fabrics:5}
(a) स्ट्राइक-स्लिप घटक \omterm{def:b1:geological-maps:attitude}{स्ट्राइक} के अनुदिश $400\cos\ang{30} = \qty{346.4}{m}$ है, और नति-सरकन घटक
$400\sin\ang{-30} = \qty{-200}{m}$ है, यानी \omterm{def:b1:geological-maps:attitude}{नति} के नीचे \qty{200}{m}; भ्रंश का \omterm{def:b1:cross-sections:throw}{विस्थापन} $200\sin\ang{50} =$ $\qty{153.2}{m}$ है;
\omterm{def:b1:cross-sections:throw}{क्षैतिज विस्थापन} $200\cos\ang{50} =$ $\qty{128.6}{m}$ है। (b) \omterm{def:b1:geological-maps:attitude}{स्ट्राइक} $000$ है, इसलिए $\mathbf s = (0, 1, 0)$; $\mathbf u = (-\cos\ang{50}, 0, \sin\ang{50}) =$
$(-0.643, 0, 0.766)$। तब, अर्थात्
$\overrightarrow{AA'} = 400(0.866\,\mathbf s - 0.5\,\mathbf u) =
(128.6, 346.4, -153.2)$ मीटर में। उसकी दिगंश $\arctan(128.6/346.4) =
\ang{20.4}$ है, और उसकी अवनति $\arcsin(153.2/400) = \ang{22.5}$, जो $\arcsin(\sin\ang{30}\sin\ang{50})$ के भी बराबर है। (c) एक \omterm{def:b1:faults-folds-fabrics:fault-types}{तिर्यक-सरकन भ्रंश}: \omterm{def:b1:faults-folds-fabrics:fault-elements}{रेक} $0$ और \ang{-90} के बीच है, $0$ के निकट: एक सामान्य घटक वाला \omterm{def:b1:faults-folds-fabrics:fault-types}{बाएँ-पार्श्वीय} भ्रंश।
\end{solution}

\begin{solution}{exo:b1:faults-folds-fabrics:6}
(a) दो भ्रंश एक-दूसरे से \ang{60} बनाते हैं; $\sigma_1$ न्यूनकोण को समद्विभाजित करता है, इसलिए वह $050$ (और $230$) के अनुदिश क्षैतिज है; $\sigma_3$ $140$ के अनुदिश क्षैतिज है; $\sigma_2$ ऊर्ध्वाधर है, उस रेखा के अनुदिश जिस पर दोनों ऊर्ध्वाधर भ्रंश मिलते हैं। हर भ्रंश $\sigma_1$ के साथ \ang{30} बनाता है, इसलिए $\ang{45} - \phi/2 =
\ang{30}$ और $\phi = \ang{30}$। (b) दोनों भ्रंश ज़मीन को चार पच्चर में काटते हैं; जिन पच्चरों में $\sigma_1$ है, वे $050$ और $230$ की ओर एक-दूसरे की ओर चलते हैं, और जिनमें $\sigma_3$ है, वे $140$ और $320$ की ओर एक-दूसरे से दूर। $020$ \omterm{def:b1:geological-maps:attitude}{स्ट्राइक} करने वाला भ्रंश लीजिए, $050$ की ओर के पच्चर (उसके पूर्व में) और $320$ की ओर के पच्चर (उसके पश्चिम में) के बीच: पश्चिमी पच्चर के सापेक्ष, पूर्वी पच्चर $230$ की ओर और $320$ से दूर चलता है, और दोनों का भ्रंश के अनुदिश $200$ की ओर घटक है। पश्चिमी पच्चर पर पूर्व की ओर मुख करके खड़ा प्रेक्षक पूर्वी पच्चर को $200$ की ओर, दाईं ओर, खिसकते देखता है: $020$ \omterm{def:b1:geological-maps:attitude}{स्ट्राइक} करने वाला भ्रंश \omterm{def:b1:faults-folds-fabrics:fault-types}{दाएँ-पार्श्वीय} है। वही तर्क, या $\sigma_1$ के बारे में सममिति, $080$ \omterm{def:b1:geological-maps:attitude}{स्ट्राइक} करने वाले भ्रंश को \omterm{def:b1:faults-folds-fabrics:fault-types}{बाएँ-पार्श्वीय} बनाती है।
\end{solution}

\begin{solution}{exo:b1:faults-folds-fabrics:7}
ऊपर की ओर के अभिलंब हैं
\begin{align*}
  \mathbf n_1 &= (\sin\ang{40}\sin\ang{120}, \sin\ang{40}\cos\ang{120}, \cos\ang{40}) = (0.557, -0.321, 0.766),\\
  \mathbf n_2 &= (\sin\ang{30}\sin\ang{260}, \sin\ang{30}\cos\ang{260}, \cos\ang{30}) = (-0.492, -0.087, 0.866).
\end{align*}
उनका सदिश-गुणनफल है, अर्थात्
$(-0.211, -0.859, -0.207)$, जो नीचे की ओर इंगित करता है: \omterm{def:b1:faults-folds-fabrics:fold-elements}{मोड़-रेखा} की अवनति, इस प्रकार $\arcsin 0.207/0.909 = \ang{13.1}$ है, और दिगंश, अर्थात् $\ang{180} + \arctan(0.211/0.859) =
\ang{193.8}$ है। $\mathbf n_1\cdot\mathbf n_2 = 0.417$; अभिलंबों के बीच का कोण \ang{65.3} है, और $\iota = \ang{114.7}$: एक \omterm{def:b1:faults-folds-fabrics:fold-classes}{खुला वलन}।
\end{solution}

\begin{solution}{exo:b1:faults-folds-fabrics:8}
$\mathbf n_1 = (0.321, 0.383, 0.866)$ परतबंदी के लिए, और
$\mathbf n_2 = (0.337, -0.925, 0.174)$ \omterm{def:b1:faults-folds-fabrics:foliation}{विदलन} के लिए। उनका सदिश-गुणनफल, नीचे की ओर लिया गया, एक ऐसी रेखा देता है जो \ang{25.4} अवनत है और \ang{074.8} की ओर है। \cref{prop:b1:faults-folds-fabrics:intersection} से यही \omterm{def:b1:faults-folds-fabrics:fold-elements}{मोड़-रेखा} की दिशा है: \omterm{def:b1:geological-maps:fold}{वलन} लगभग \ang{25} पूर्व-उत्तर-पूर्व की ओर अवनत है, जिसकी पुष्टि \omterm{def:b1:faults-folds-fabrics:fold-elements}{पार्श्व} के किसी भी अन्य उद्भेद को करनी चाहिए।
\end{solution}

\begin{solution}{exo:b1:faults-folds-fabrics:9}
(a) दक्षिण में, उत्तर की ओर मुख करके खड़ा प्रेक्षक भ्रंश के उत्तरी भाग को पूर्व की ओर, दाईं ओर, विस्थापित देखता है: भ्रंश \omterm{def:b1:faults-folds-fabrics:fault-types}{दाएँ-पार्श्वीय} है। (b) डाइक भ्रंश के लंबवत एक ऊर्ध्वाधर तल है, इसलिए \omterm{def:b1:geological-maps:attitude}{स्ट्राइक} के अनुदिश उसका \omterm{def:b1:cross-sections:throw}{विस्थापन} स्ट्राइक-स्लिप घटक है: \qty{350}{m}, जो पूरा सरकन है यदि \omterm{def:b1:faults-folds-fabrics:fault-rocks}{स्लिकेनलाइन} क्षैतिज हों। (c) तब सरकन ऊर्ध्वाधर भ्रंश-तल में \omterm{def:b1:geological-maps:attitude}{स्ट्राइक} से \ang{20} बनाता है: उसका क्षैतिज घटक अब भी \qty{350}{m} है, इसलिए \omterm{def:b1:faults-folds-fabrics:fault-elements}{निवल सरकन} $350/\cos\ang{20} = \qty{372.5}{m}$ है और उसका ऊर्ध्वाधर घटक $350\tan\ang{20} = \qty{127.4}{m}$। उत्तरी खंड पूर्व की ओर गया, और \omterm{def:b1:faults-folds-fabrics:fault-rocks}{स्लिकेनलाइन} पूर्व की ओर नीचे उतरती हैं, इसलिए वह भी \qty{127.4}{m} नीचे गया। एक ऊर्ध्वाधर डाइक वह ऊर्ध्वाधर गति नहीं दिखा सकती, क्योंकि उसे उसके अपने अनुदिश खिसकाने से उसका निशान नहीं बदलता।
\end{solution}

\begin{solution}{exo:b1:faults-folds-fabrics:10}
(a) भ्रंश का अभिलंब $\ang{90} - \theta$ बनाता है $\sigma_1$ के साथ, इसलिए \cref{prop:b1:stress-strain-rheology:mohr} से,
$\sigma_n' = \sigma_1'\sin^2\theta + \sigma_3'\cos^2\theta$ और $\tau = (\sigma_1' - \sigma_3')\sin\theta\cos\theta$। सरकन तब शुरू होता है जब $\tau = \mu\sigma_n'$; $\sigma_3'\cos^2\theta$ से भाग देकर और $R = \sigma_1'/\sigma_3'$ लिखकर, $(R - 1)\tan\theta = \mu(R\tan^2\theta + 1)$, इसलिए $R\tan\theta(1 - \mu\tan\theta) = \tan\theta + \mu$, जो सूत्र है। (b) $\sigma_1$ ऊर्ध्वाधर होने पर, $\theta$ \omterm{def:b1:geological-maps:attitude}{नति} को \ang{90} में से घटाने पर मिलता है: $\theta = \ang{30}$ पर $R = 3.12$; $\theta = \ang{45}$ पर $R = 4.00$; $\theta = \ang{55}$ पर $R = 9.92$। (c) हर भाजक $\tan\theta = 1/\mu$, $\theta = \ang{59.0}$ पर शून्य है; उससे आगे $R$ ऋणात्मक होता, अर्थात धनात्मक प्रभावी प्रतिबलों का कोई अनुपात भ्रंश को नहीं सरका सकता, चाहे $\sigma_1'$ कितना भी बड़ा हो। ऊर्ध्वाधर $\sigma_1$ के अंतर्गत, $\mu = 0.6$ के साथ, \ang{31.0} से कम नत \omterm{def:b1:faults-folds-fabrics:fault-types}{सामान्य भ्रंश} कभी नहीं सरक सकते। (d) सबसे उपयुक्त दिशा वाला तल $\ang{45} - \phi/2$ पर है, $\sigma_1$ से। $\mu = 0.6$, $\phi = \ang{31.0}$ के लिए: \omterm{def:b1:faults-folds-fabrics:fault-types}{सामान्य भ्रंश} \ang{60.5} नत हैं, अधिक्षेप \ang{29.5}। $\mu = 0.85$, $\phi = \ang{40.4}$ के लिए: \ang{65.2} और \ang{24.8}।
\end{solution}

\begin{solution}{exo:b1:faults-folds-fabrics:11}
उत्तर की ओर देखने पर पूर्व दाईं ओर है, और पश्चिम की ओर नत तल दाईं ओर उठते हैं। (a) B पर परतबंदी \omterm{def:b1:faults-folds-fabrics:foliation}{विदलन} की ही ओर नत है, अधिक तीव्रता से: \cref{prop:b1:faults-folds-fabrics:cleavage-bedding}(a) से, B एक उलटे \omterm{def:b1:faults-folds-fabrics:fold-elements}{पार्श्व} पर है; A और C, जहाँ परतबंदी कम खड़ी है, सीधे हैं। (b) \omterm{def:b1:faults-folds-fabrics:foliation}{विदलन} \ang{60} पर उठता है; A और C पर परतबंदी \ang{40} पर उठती है, और एक से दूसरे का न्यूनकोणीय घूर्णन दक्षिणावर्त है, \ang{20} से होकर, जो \omterm{def:b1:faults-folds-fabrics:foliation}{विदलन} के शीर्ष को पूर्व की ओर ले जाता है; B पर परतबंदी \ang{80} पर उठती है और घूर्णन वामावर्त है, \ang{20} से होकर, जो उसे पश्चिम की ओर ले जाता है। (c) संवेदना A और B के बीच बदलती है और फिर B और C के बीच: हर अंतराल में एक मोड़ है। A पर, जो सीधे \omterm{def:b1:faults-folds-fabrics:fold-elements}{पार्श्व} पर है, \omterm{def:b1:faults-folds-fabrics:foliation}{विदलन} का शीर्ष पूर्व की ओर, \omterm{def:b1:faults-folds-fabrics:fold-elements}{पार्श्व} के शीर्ष पर \omterm{def:b1:geological-maps:fold}{अपनति} की ओर जाता है: A और B के बीच का मोड़ एक \omterm{def:b1:geological-maps:fold}{अपनति} है, और B और C के बीच का एक \omterm{def:b1:geological-maps:fold}{अभिनति}। (d) प्रमेय के प्रमाण के संकेतन में, सीधे \omterm{def:b1:faults-folds-fabrics:fold-elements}{पार्श्व} $\gamma - \iota/2 = \ang{40}$ नत हैं और उलटे \omterm{def:b1:faults-folds-fabrics:fold-elements}{पार्श्व} $\gamma + \iota/2 =
\ang{80}$, इसलिए $\gamma = \ang{60}$, जैसा \omterm{def:b1:faults-folds-fabrics:foliation}{विदलन} कहता है, और $\iota =
\ang{40}$: \omterm{def:b1:faults-folds-fabrics:fold-classes}{सघन वलन}। \omterm{def:b1:faults-folds-fabrics:fold-elements}{अक्षीय तल} उत्तर की ओर \omterm{def:b1:geological-maps:attitude}{स्ट्राइक} करते हैं और \ang{60} पश्चिम की ओर नत हैं, विन्यास $270/60$। छोटे \omterm{def:b1:faults-folds-fabrics:fold-elements}{पार्श्व} ऊर्ध्वाधर से आगे मुड़ गए हैं, उनके शीर्ष पूर्व की ओर: \omterm{def:b1:geological-maps:fold}{वलन} पूर्व की ओर झुकते हैं, जैसा अध्याय के चित्र में है।
\end{solution}

\begin{solution}{exo:b1:faults-folds-fabrics:12}
(a) एक सममित शेवरॉन \omterm{def:b1:geological-maps:fold}{वलन} में हर \omterm{def:b1:faults-folds-fabrics:fold-elements}{पार्श्व} \omterm{def:b1:faults-folds-fabrics:fold-elements}{अक्षीय तल} के साथ $\iota/2$ कोण बनाता है, जो संकुचन की दिशा के लंबवत है; $l$ लंबाई का एक \omterm{def:b1:faults-folds-fabrics:fold-elements}{पार्श्व} वलनों के आर-पार $l\sin(\iota/2)$ फैलता है, और पूरी परत $L_0\sin(\iota/2)$। (b) $\sin\ang{30} = 0.5$: \qty{50}{\percent} संकुचन। (c) $\sin(\iota/2) = 0.7$ से $\iota = \ang{88.9}$। (d) $\iota = \ang{156}$ और $\sin\ang{78} = 0.978$: \qty{2.2}{\percent} संकुचन, जो $1 - \cos\ang{12}$ के भी बराबर है।
\end{solution}

\begin{solution}{pb:b1:faults-folds-fabrics:1}
\textbf{1.} $\mathbf n = (\sin\delta\sin\alpha, \sin\delta\cos\alpha,
\cos\delta)$ से
\[
  \mathbf n_1 = (0.20855, 0.03415, 0.97742), \qquad
  \mathbf n_2 = (-0.20855, 0.03415, 0.97742).
\]

\textbf{2.} $\mathbf n_1\cdot\mathbf n_2 = -0.04349 + 0.00117 + 0.95535 =
0.91302$, और अभिलंबों के बीच का कोण \ang{24.1} है।

\textbf{3.} \cref{prop:b1:faults-folds-fabrics:hinge} से,
$\iota = \ang{180} - \ang{24.1} = \ang{155.9}$।

\textbf{4.} \ang{170} और \ang{90} के बीच: एक \omterm{def:b1:faults-folds-fabrics:fold-classes}{खुला वलन}।

\textbf{5.} \omterm{def:b1:faults-folds-fabrics:fold-elements}{पार्श्व} $080.7$ और $279.3$ की ओर नत हैं, ठीक पूर्व और पश्चिम की ओर नहीं,
क्योंकि \omterm{def:b1:geological-maps:fold}{वलन} अवनत है: उनकी सबसे खड़ी रेखाएँ परिच्छेद में नहीं हैं, जो केवल एक \omterm{def:b1:cross-sections:apparent-dip}{आभासी नति} दिखाता है। पूर्व–पश्चिम काट में,
\cref{thm:b1:cross-sections:apparent-dip} से, $\tan\beta = \tan\ang{12.2}
\cos\ang{9.3} = 0.2134$, तो $\beta = \ang{12.0}$, और $\ang{180} -
2\times\ang{12.0} = \ang{155.9}$ (काट से प्रोफ़ाइल \ang{2} झुकी है, परंतु इस परिशुद्धता पर इससे कुछ नहीं बदलता)।

\textbf{6.} उसके \omterm{def:b1:faults-folds-fabrics:fold-elements}{पार्श्व} \omterm{def:b1:geological-maps:fold}{अपनति} के पार्श्वों के समांतर हैं, इसलिए उसका \omterm{def:b1:faults-folds-fabrics:fold-elements}{अंतर-पार्श्व कोण} भी
\ang{155.9} है।

\textbf{7.} अब $\mathbf n_1\times\mathbf n_2 = (0.03415\times0.97742 -
0.97742\times0.03415,\; 0.97742\times(-0.20855) - 0.20855\times0.97742,\;
0.20855\times0.03415 + 0.03415\times0.20855) = (0, -0.40768, 0.01424)$।
नीचे की ओर लेने पर हमें $(0, 0.40768, -0.01424)$: \omterm{def:b1:faults-folds-fabrics:fold-elements}{मोड़-रेखा} की दिगंश $000$ है और वह
उत्तर की ओर $\arctan(0.01424/0.40768) = \ang{2.0}$ अवनत है।

\textbf{8.} \omterm{def:b1:faults-folds-fabrics:fold-elements}{अक्षीय तल} \omterm{def:b1:faults-folds-fabrics:fold-elements}{मोड़-रेखा} को समाहित करता है और दोनों पार्श्वों को समद्विभाजित करता है;
उसका अभिलंब $\mathbf n_1 - \mathbf n_2 = (0.41710, 0, 0)$ के अनुदिश है, क्षैतिज और
पूर्व–पश्चिम। वह ऊर्ध्वाधर है और उत्तर–दक्षिण \omterm{def:b1:geological-maps:attitude}{स्ट्राइक} करता है: $090/90$। \omterm{def:b1:geological-maps:fold}{अपनति}
सीधी है।

\textbf{9.} एक खुला, \omterm{def:b1:faults-folds-fabrics:fold-classes}{सीधा वलन}, जो उत्तर की ओर \ang{2.0} अवनत है।

\textbf{10.} \cref{prop:b1:geological-maps:folds} से, एक \omterm{def:b1:geological-maps:fold}{अपनति} के उद्भेद अवनति की दिशा में बंद होते हैं: उत्तर की ओर। उसके केंद्र में सबसे पुरानी चट्टानें हैं: आवरण के आधार पर बलुआ पत्थर और, जहाँ नदी ने सबसे गहरा काटा है, उसके नीचे शिस्ट और ग्रेनाइट का आधार।

\textbf{11.} अक्षीय-तलीय होने से, \omterm{def:b1:faults-folds-fabrics:foliation}{विदलन} ऊर्ध्वाधर है और उत्तर–दक्षिण \omterm{def:b1:geological-maps:attitude}{स्ट्राइक} करता है,
$090/90$। \cref{prop:b1:faults-folds-fabrics:intersection} से, उसका परतबंदी के साथ
\omterm{def:b1:faults-folds-fabrics:lineation}{प्रतिच्छेद-रेखन} \omterm{def:b1:faults-folds-fabrics:fold-elements}{मोड़-रेखा} के समांतर है: दिगंश
$000$, अवनति \ang{2.0}। सीधे: तल $080.7/12.2$ ऊर्ध्वाधर तल से मिलता है
जो उत्तर की ओर \omterm{def:b1:geological-maps:attitude}{स्ट्राइक} करता है, एक ऐसी रेखा के साथ जिसकी अवनति $\arctan(\tan\ang{12.2}
\cos\ang{80.7}) = \ang{2.0}$ है, $000$ की ओर।

\textbf{12.} उत्तर की ओर देखने पर पूर्वी \omterm{def:b1:faults-folds-fabrics:fold-elements}{पार्श्व} बाईं ओर \ang{12.0} उठता है
और \omterm{def:b1:faults-folds-fabrics:foliation}{विदलन} ऊर्ध्वाधर है: \omterm{def:b1:faults-folds-fabrics:foliation}{विदलन} से परतबंदी तक का न्यूनकोणीय घूर्णन, \ang{78.0} से होकर, वामावर्त है और \omterm{def:b1:faults-folds-fabrics:foliation}{विदलन} के शीर्ष को पश्चिम की ओर ले जाता है। \omterm{def:b1:faults-folds-fabrics:fold-elements}{पार्श्व} सीधा है, इसलिए
\cref{prop:b1:faults-folds-fabrics:cleavage-bedding}(b) से \omterm{def:b1:geological-maps:fold}{अपनति} का मोड़ पश्चिम में है, $x = \qty{3.5}{km}$ पर, जैसा मानचित्र दिखाता है।

\textbf{13.} \omterm{def:b1:geological-maps:attitude}{स्ट्राइक} $120 - 90 = 030$ है; तल $120$ की ओर नत है, इसलिए
\omterm{def:b1:faults-folds-fabrics:fault-elements}{लटकती दीवार} उस रेखा के पूर्व-दक्षिण-पूर्व की ओर है।

\textbf{14.} भ्रंश सामान्य है: $\sigma_1$ ऊर्ध्वाधर था।

\textbf{15.} $\sigma_2$ \omterm{def:b1:geological-maps:attitude}{स्ट्राइक} के अनुदिश क्षैतिज है, $030$--$210$;
$\sigma_3$ उसके आर-पार क्षैतिज है, $120$--$300$।

\textbf{16.} $\ang{45} + \phi/2 = \ang{60}$ से $\phi = \ang{30}$ मिलता है, और
$\mu = \tan\ang{30} = 0.577$।

\textbf{17.} संयुग्मी समुच्चय भी $030$ \omterm{def:b1:geological-maps:attitude}{स्ट्राइक} करता है और दूसरी ओर \ang{60} नत है:
$300/60$। दोनों समुच्चय $\ang{180} - 2\times\ang{60} = \ang{60}$ पर मिलते हैं, और
ऊर्ध्वाधर $\sigma_1$ न्यूनकोण को समद्विभाजित करता है।

\textbf{18.} $\mu = 0.85$, $\phi = \ang{40.4}$ के साथ सबसे उपयुक्त दिशा वाले \omterm{def:b1:faults-folds-fabrics:fault-types}{सामान्य
भ्रंश} \ang{65.2} नत हैं; वार्डन भ्रंश, \ang{60} पर, \ang{5.2} दूर है।
\cref{exo:b1:faults-folds-fabrics:10} से, वह अनुपात $\sigma_1'/\sigma_3' = (1 + 0.85\cot\ang{30})/(1 - 0.85\tan\ang{30}) = 4.85$
पर सरकता है, जबकि सबसे उपयुक्त तल के लिए $4.68$ है: वह लगभग, पर ठीक नहीं, सबसे उपयुक्त दिशा में है, और उसी शासन में फिर सरकने के लिए केवल \qty{4}{\percent} अधिक प्रतिबल अनुपात चाहिए।

\textbf{19.} सरकन \omterm{def:b1:geological-maps:attitude}{नति} के नीचे की ओर है: $r = \ang{-90}$; \omterm{def:b1:faults-folds-fabrics:fault-elements}{सरकन-सदिश}
$120$ की दिशा में है और \ang{60} अवनत है।

\textbf{20.} $(\cos\ang{60}\sin\ang{120}, \cos\ang{60}\cos\ang{120},
-\sin\ang{60}) = (0.4330, -0.2500, -0.8660)$।

\textbf{21.} $e = 200/\tan\ang{60} = \qty{115.5}{m}$।

\textbf{22.} पूर्व: $115.47\sin\ang{120} = \qty{100.0}{m}$; उत्तर:
$115.47\cos\ang{120} = \qty{-57.7}{m}$, यानी दक्षिण की ओर \qty{57.7}{m}।

\textbf{23.} $L = 200/\sin\ang{60} = \qty{230.9}{m}$, और
$200^2 + 115.47^2 = 40\,000 + 13\,333 = 53\,333 = 230.94^2$। भ्रंश का \omterm{def:b1:cross-sections:throw}{विस्थापन} और
\omterm{def:b1:cross-sections:throw}{क्षैतिज विस्थापन} सरकन के ऊर्ध्वाधर और क्षैतिज घटक हैं, जो \omterm{def:b1:geological-maps:attitude}{नति} के ऊर्ध्वाधर तल में
है: वे एक समकोण त्रिभुज की भुजाएँ हैं जिसका कर्ण सरकन है।

\textbf{24.} निकटतम मीटर तक, वार्डन भ्रंश पर \omterm{def:b1:faults-folds-fabrics:fault-elements}{निवल सरकन} है
\textbf{\qty[text-series-to-math=true]{231}{m}}।
\end{solution}
